% ===== PRÉAMBULE CANONIQUE — à coller VERBATIM en tête de CHAQUE .tex =====
\documentclass[11pt,a4paper]{article}
\usepackage[utf8]{inputenc}\usepackage[T1]{fontenc}\usepackage{lmodern}
\usepackage[french]{babel}
\usepackage{amsmath,amssymb,mathtools}\usepackage{icomma}
\usepackage[version=4]{mhchem}          % \ce{...} : équations & formules
\usepackage{siunitx}\sisetup{locale=FR} % \qty{}{}, \si{}, \ang{}
\usepackage{chemfig}                    % \chemfig{...} : structures développées
\usepackage{xcolor}\usepackage{colortbl}
\usepackage{tikz}\usepackage{pgfplots}\pgfplotsset{compat=1.18}
\usepackage[most]{tcolorbox}\usepackage{enumitem}\usepackage{array,booktabs}
\usepackage[a4paper,margin=2.2cm]{geometry}
\usetikzlibrary{arrows.meta,calc,angles,quotes,positioning,patterns,backgrounds,shapes.geometric}
\definecolor{primary}{RGB}{222,49,99}\definecolor{primarydark}{RGB}{150,22,55}
\definecolor{defcol}{RGB}{0,114,178}\definecolor{methcol}{RGB}{52,92,148}
\definecolor{excol}{RGB}{0,135,90}\definecolor{attncol}{RGB}{200,40,55}
\tcbset{boxbase/.style={enhanced,breakable,boxrule=0.9pt,arc=2.5pt,
  left=3mm,right=3mm,top=2.3mm,bottom=2.3mm,fonttitle=\bfseries,coltitle=white}}
% --- boîtes TUTORIEL (noms EXACTS, ne pas renommer) ---
\newtcolorbox{cle}[1][]{boxbase,colback=defcol!6,colframe=defcol,
  title={\textsc{À retenir}\ifx\relax#1\relax\relax\else\textmd{ : \itshape #1}\fi}}
\newtcolorbox{methode}[1][]{boxbase,colback=methcol!7,colframe=methcol,sharp corners,
  title={\textsc{Méthode}\ifx\relax#1\relax\relax\else\textmd{ --- \itshape #1}\fi}}
\newtcolorbox{exemplebox}[1][]{boxbase,colback=excol!5,colframe=excol,
  title={\textsc{Exemple}\ifx\relax#1\relax\relax\else\textmd{ : \itshape #1}\fi}}
\newtcolorbox{piege}[1][]{boxbase,colback=attncol!6,colframe=attncol,
  title={\textsc{Piège}\ifx\relax#1\relax\relax\else\textmd{ : \itshape #1}\fi}}
\newtcolorbox{remarque}[1][]{boxbase,colback=black!4,colframe=black!45,
  title={\textsc{Remarque}\ifx\relax#1\relax\relax\else\textmd{ : \itshape #1}\fi}}
% --- boîtes EXERCICE / CORRIGÉ (noms EXACTS, auto-comptées) ---
\newtcolorbox[auto counter]{exo}[1][]{boxbase,colback=primary!4,colframe=primary,
  title={\textsc{Exercice}~\thetcbcounter\ifx\relax#1\relax\relax\else\textmd{ --- \itshape #1}\fi}}
\newtcolorbox[auto counter]{corrige}[1][]{boxbase,colback=excol!4,colframe=excol,
  title={\textsc{Corrigé}~\thetcbcounter\ifx\relax#1\relax\relax\else\textmd{ --- \itshape #1}\fi}}
\newcommand{\R}{\mathbb{R}}\newcommand{\N}{\mathbb{N}}\newcommand{\Z}{\mathbb{Z}}\newcommand{\ds}{\displaystyle}
% (un \usepackage{fancyhdr}+en-tête est possible ; il est ignoré côté web)
% =========================================================================
\begin{document}
% THEME_TITLE: Nombre d'oxydation et décompte électronique

\begin{cle}[à quoi sert le nombre d'oxydation]
Le \textbf{nombre d'oxydation} (n.o., noté en chiffres romains : $+\mathrm{I}$, $-\mathrm{II}$\ldots) d'un
élément dans une espèce est la charge \emph{fictive} qu'il porterait si chaque liaison était attribuée à
l'atome le plus électronégatif. Il est la clé de toute la nomenclature : c'est lui qui fixe la
\textbf{valence} d'un métal à plusieurs états (fer(II) vs fer(III)) et qui permet de vérifier une formule.
Pour un ion monoatomique, le n.o.\ est \emph{égal à la charge réelle}.
\end{cle}

\begin{methode}[déterminer le n.o.\ d'un élément — les règles dans l'ordre]
On applique les conventions suivantes, puis on résout l'équation de la charge :
\begin{enumerate}[label=\arabic*.,leftmargin=2.0em,topsep=2pt,itemsep=1pt]
  \item \textbf{corps simple} (élément seul : \ce{O2}, \ce{Fe}, \ce{Cl2}) $\Rightarrow$ n.o.\ $=0$ ;
  \item \textbf{ion monoatomique} $\Rightarrow$ n.o.\ $=$ charge (\ce{Cl-} : $-\mathrm{I}$ ; \ce{Ca^2+} : $+\mathrm{II}$) ;
  \item \textbf{oxygène} $\Rightarrow -\mathrm{II}$ \quad(\emph{exceptions} : peroxydes $-\mathrm{I}$) ;
  \item \textbf{hydrogène} $\Rightarrow +\mathrm{I}$ \quad(\emph{exceptions} : hydrures métalliques $-\mathrm{I}$) ;
  \item alcalins (\ce{Na},\ce{K}) $+\mathrm{I}$ ; alcalino-terreux (\ce{Ca},\ce{Mg}) $+\mathrm{II}$ ; fluor $-\mathrm{I}$ ;
  \item \textbf{équation de fermeture} : la \emph{somme} de tous les n.o.\ (chacun multiplié par son indice)
        égale la \textbf{charge globale} de l'espèce (0 pour une molécule neutre).
\end{enumerate}
On note $x$ le n.o.\ inconnu, on écrit l'équation et on isole $x$.
\end{methode}

\begin{exemplebox}[soufre dans \ce{SO4^2-} et chrome dans \ce{Cr2O7^2-}]
\textbf{\ce{SO4^2-} :} l'oxygène vaut $-\mathrm{II}$, la charge globale est $-2$. Avec $x$ le n.o.\ du soufre :
\[ x + 4\times(-\mathrm{II}) = -2 \;\Rightarrow\; x - 8 = -2 \;\Rightarrow\; \boxed{x = +\mathrm{VI}}. \]
\textbf{\ce{Cr2O7^2-} :} il y a \emph{deux} atomes de chrome et sept d'oxygène :
\[ 2x + 7\times(-\mathrm{II}) = -2 \;\Rightarrow\; 2x = +12 \;\Rightarrow\; \boxed{x = +\mathrm{VI}}. \]
\end{exemplebox}

\begin{cle}[du n.o.\ au nombre d'électrons d'un ion]
Un atome neutre de numéro atomique $Z$ possède \emph{$Z$ électrons}. Un ion de charge $q$ (égale au n.o.\
pour un ion monoatomique) a \emph{gagné} ou \emph{perdu} des électrons :
\[ \boxed{\ N_{\text{électrons}} = Z - q\ } \qquad (q>0 \text{ : cation, a perdu ; } q<0 \text{ : anion, a gagné}). \]
Ainsi \ce{Cu+} ($Z=29$) possède $29-1=28$ électrons ; \ce{O^2-} ($Z=8$) en possède $8-(-2)=10$.
\end{cle}

\begin{piege}[les trois erreurs classiques]
\begin{itemize}[leftmargin=1.4em,topsep=2pt,itemsep=1pt]
  \item un \textbf{métal} (cuivre, fer, sodium\ldots) ne prend \emph{jamais} de n.o.\ négatif à l'état
        naturel : on écarte d'office « cuivre $-\mathrm{I}$ ».
  \item ne pas oublier l'\textbf{indice} : dans \ce{Cr2O7^2-} il y a $2$ chromes, donc c'est $2x$, pas $x$.
  \item le n.o.\ de \ce{O} est $-\mathrm{II}$ \emph{sauf} dans les peroxydes (\ce{H2O2}, \ce{Na2O2} : $-\mathrm{I}$),
        et celui de \ce{H} est $+\mathrm{I}$ \emph{sauf} dans les hydrures (\ce{NaH}, \ce{CaH2} : $-\mathrm{I}$).
\end{itemize}
\end{piege}

\begin{remarque}[pourquoi c'est utile en nomenclature]
Le n.o.\ tranche les cas de \textbf{valence variable} : dans \ce{FeO}, le fer est $+\mathrm{II}$
(oxyde de fer(II)) ; dans \ce{Fe2O3}, il est $+\mathrm{III}$ (oxyde de fer(III)). Même couple d'éléments,
mais un n.o.\ différent \emph{impose} une formule et un nom différents.
\end{remarque}

\end{document}
